\documentclass{article}

\usepackage{amsmath}
\usepackage{amssymb}

\title{Data Compression using Singular Value Decomposition (SVD)}
\author{Mardhav Rayapati}
\date{}

\begin{document}

\maketitle
\newpage

\section{Singular Value Decomposition}

\subsection{Introduction}

Singular Value Decomposition (SVD) is a matrix factorization technique
that decomposes a matrix into three matrices. For a matrix
$A \in \mathbb{R}^{m \times n}$, its singular value decomposition is

\[
    A = U \Sigma V^{T},
\]

where

\[
    U \in \mathbb{R}^{m \times m},
    \qquad
    \Sigma \in \mathbb{R}^{m \times n},
    \qquad
    V \in \mathbb{R}^{n \times n}.
\]

The matrices $U$ and $V$ are orthogonal matrices, satisfying

\[
    U^{T}U = I_m,
    \qquad
    V^{T}V = I_n,
\]

while $\Sigma$ is a diagonal matrix containing the singular values
of $A$.

The singular values are conventionally arranged in non-increasing
order:

\[
    \sigma_1 \geq \sigma_2 \geq \cdots \geq
    \sigma_r \geq 0,
\]

where

\[
    r = \operatorname{rank}(A).
\]

Equivalently, the decomposition can be written as

\[
    A =
    \sum_{i=1}^{r}
    \sigma_i u_i v_i^{T},
\]

where $u_i$ and $v_i$ are the corresponding left and right singular
vectors.

\subsection{Connection to Eigenvalue Decomposition}

The singular values of $A$ are related to the eigenvalues of
$A^{T}A$.

Since

\[
    A = U\Sigma V^{T},
\]

we obtain

\[
    A^{T}A
    =
    V\Sigma^{T}U^{T}U\Sigma V^{T}.
\]

Using the orthogonality of $U$,

\[
    U^{T}U = I,
\]

giving

\[
    A^{T}A = V\Sigma^{T}\Sigma V^{T}.
\]

Therefore, $V$ contains the eigenvectors of $A^{T}A$, while the
eigenvalues are the squared singular values:

\[
    \lambda_i = \sigma_i^2.
\]

Consequently,

\[
    \sigma_i = \sqrt{\lambda_i}.
\]

This relationship provides one way of computing the SVD and forms
the basis for the eigenvalue-based approach explored in this project.

\section{Step 1: Householder Reflector Matrix}

In the eigenvalue-based computation of the SVD, the symmetric
matrix

\[
    B = A^T A
\]

is first constructed.\\

Geometrically, a Householder reflection matrix reflects a vector
across a hyperplane perpendicular to a given vector.\\

For a nonzero vector $u$, a Householder reflector is generally
defined as

\[
    H = I - 2\frac{uu^T}{u^T u}.
\]

If $u$ is a unit vector, then $u^T u = 1$, and the expression
simplifies to

\[
    H = I - 2uu^T.
\]

Every Householder reflector has the following properties:

\begin{itemize}

    \item Symmetric:
          \[
              H^T = H
          \]

    \item Orthogonal:
          \[
              H^T H = I
          \]

\end{itemize}

\section{Step 2: Tridiagonalization}

A symmetric matrix $B$ can be transformed into a symmetric
tridiagonal matrix by applying a series of Householder
transformations.\\

Each transformation is an orthogonal similarity transformation
of the form

\[
    B_{k+1} = H_k B_k H_k^T,
\]

where $H_k$ is a Householder reflector.\\

After applying the required sequence of transformations, the
matrix is reduced to a symmetric tridiagonal matrix of the form

\[
    T =
    \begin{bmatrix}
        b_{11} & b_{12} & 0      & 0      & 0      & \dots \\
        b_{12} & b_{22} & b_{23} & 0      & 0      & \dots \\
        0      & b_{23} & b_{33} & b_{34} & 0      & \dots \\
        0      & 0      & b_{34} & b_{44} & b_{45} & \dots \\
        \vdots & \vdots & \vdots & \vdots & \vdots &
    \end{bmatrix}.
\]

Since each Householder transformation is orthogonal, the
transformation preserves the eigenvalues of the matrix. Therefore,

\[
    \lambda_i(T) = \lambda_i(B).
\]

\section{Step 3: QR Iteration}

The tridiagonal matrix $T$ obtained from the previous step is
used to compute its eigenvalues. This is done using the QR
iteration algorithm.

Let

\[
    T_0 = T.
\]

At each iteration, the current matrix $T_k$ is factorized using
the QR decomposition:

\[
    T_k = Q_k R_k,
\]

where $Q_k$ is orthogonal and $R_k$ is upper triangular.\\

The next iterate is then obtained by reversing the order of the
factors:

\[
    T_{k+1} = R_k Q_k.
\]

Since

\[
    T_k = Q_k R_k,
\]

we have

\[
    T_{k+1}
    = R_k Q_k
    = Q_k^T T_k Q_k.
\]

Therefore, $T_{k+1}$ is orthogonally similar to $T_k$ and has
the same eigenvalues.\\

Repeating this process,

\[
    T_0 \rightarrow T_1 \rightarrow T_2
    \rightarrow \cdots \rightarrow T_k,
\]

causes the off-diagonal elements to converge towards zero. Thus,

\[
    T_k \rightarrow D,
\]

where $D$ is a diagonal matrix:

\[
    D =
    \begin{bmatrix}
        \lambda_1 & 0         & \cdots & 0         \\
        0         & \lambda_2 & \cdots & 0         \\
        \vdots    & \vdots    & \ddots & \vdots    \\
        0         & 0         & \cdots & \lambda_n
    \end{bmatrix}.
\]

The diagonal entries of $D$ are therefore the eigenvalues of
the original matrix $B = A^T A$.\\

The singular values of $A$ can then be obtained from these
eigenvalues:

\[
    \sigma_i = \sqrt{\lambda_i}.
\]

\section{Step 3: Construct the Eigenvectors of B using the Eigenvalues}

\end{document}